% Bewertungspaket zum Gesamttest «Lineare Gleichungen» (GF 2.2a) — Quelle des PDFs (python3 scripts/build-lp-pdf.py /lineare-gleichungen/).
% Ganze Punkte je Teilschritt; Werte mit python3 nachgerechnet (11.10.2026, Fehlerwerte eingeschlossen).
\documentclass[11pt]{article}
\usepackage{../lp-druck}
\renewcommand{\lpname}{Lineare Gleichungen}
\renewcommand{\lpbereich}{Grundlagenfach}
\renewcommand{\lpteilgebiet}{2.2a}
\renewcommand{\lpfuss}{Bewertungspaket}
\newcolumntype{P}{>{\centering\arraybackslash}p{10mm}}
\newenvironment{raster}{\par\smallskip\noindent\tabularx{\linewidth}{@{}>{\raggedright\arraybackslash}p{38mm}P X@{}}\toprule
  \small\textsc{Teilschritt} & \small\textsc{P} & \small\textsc{Musterlösung}\\\midrule}{\bottomrule\endtabularx\par}
\newcommand{\fehler}[1]{\par\smallskip{\small\textcolor{lprot}{\bfseries Typische Fehler:} #1}}
\newcommand{\E}{\,{\footnotesize\bfseries\color{lpgruen}(E)}}
\newcommand{\A}{\,{\footnotesize\bfseries\color{lpgruen}(A)}}
\begin{document}
\lpkopf{Bewertungspaket zum Gesamttest}{}

\begin{lpachtung}{Erst lösen, dann öffnen}
Dieses Paket enthält die Musterlösung. Wer es vor dem Lösen liest, misst am Ende nur, wie gut das Abschreiben gelingt.
\end{lpachtung}

\section*{1 · So gehst du vor}
\begin{enumerate}[leftmargin=*,itemsep=2pt]
\item \textbf{Gesamttest lösen} — vollständig, mit Rechenweg, ohne dieses Paket und ohne Taschenrechner.
\item \textbf{Fotografieren oder scannen}, gut lesbar, eine Seite pro Bild. \textbf{Kein Name} auf den Bildern.
  Handyfotos tragen oft unsichtbar Standort, Zeit und Gerät mit (Metadaten). Lade darum lieber einen Scan oder ein
  Bildschirmfoto des Fotos hoch, oder entferne vorher die Standortangaben.
\item \textbf{Bewerten lassen.} Ob du dafür eine KI nutzt und welche, entscheidest du selbst und in eigener Verantwortung —
  je nachdem, was dir aufgrund deines Alters und deiner persönlichen Situation erlaubt ist (Altersgrenzen und
  Nutzungsbedingungen des Anbieters, allenfalls Einverständnis deiner Eltern). Lade nur deine Lösung und dieses PDF hoch,
  keine persönlichen Daten, und füge den Auftrag aus Abschnitt~2 ein. Ohne KI kannst du dich mit Abschnitt~3 selbst bewerten.
\item \textbf{Rückmeldung prüfen.} Stimmt, was die KI aus deiner Handschrift gelesen hat? Eine KI kann sich irren —
  im Zweifel gilt das Raster in Abschnitt~3.
\item \textbf{Gezielt wiederholen} nach Abschnitt~4.
\end{enumerate}

\section*{2 · Auftrag an die KI}
\begin{tcolorbox}[colback=lplinie!25,colframe=lplinie,boxrule=0.5pt,arc=1mm,fontupper=\small]
Du bist Korrektorin oder Korrektor für Mathematik an einer Schweizer Berufsmaturitätsschule. Im Anhang findest du (1)~das Bewertungspaket zum Gesamttest «Lineare Gleichungen» mit Musterlösung und Punkteraster und (2)~Fotos meiner handschriftlichen Lösung.\par\medskip
Geh so vor:\par
1. Schreib zu jeder Aufgabe G1 bis G7 zuerst ab, was du in meiner Lösung liest — Rechenweg und Ergebnis. Was du nicht sicher lesen kannst, markierst du mit [unsicher]. Nicht raten.\par
2. Vergib die Punkte genau nach dem Raster im Bewertungspaket (Abschnitt 3), ganze Punkte je Zeile. Ziehe einen Fehler nur einmal ab: Rechne ich mit einem falschen Zwischenergebnis richtig weiter, gibt es die Folgepunkte. Ausnahme: Zeilen mit \textbf{(E)} gibt es nur für das richtige Ergebnis, nie als Folgepunkt. Die Zeilen «Typische Fehler» sagen, welche Rasterzeile bei diesem Fehler 0 Punkte gibt — sie sind kein zusätzlicher Abzug.\par
3. Andere richtige Lösungswege zählen voll. Gleichwertige Schreibweisen zählen voll: Dezimalpunkt oder Dezimalkomma, Bruch oder Dezimalzahl ($\tfrac75 = 1.4$), gekürzter oder ungekürzter Bruch, $\mathbb{L} = \{6\}$ oder «$x = 6$», $\mathbb{L} = \{\,\}$ oder «keine Lösung», $\mathbb{L} = \mathbb{R}$ oder «jede Zahl ist Lösung» oder «unendlich viele Lösungen», eine Normalform als Vielfaches ($-5x + 7 = 0$ oder $10x - 14 = 0$ statt $5x - 7 = 0$) — ausser wo das Raster eine bestimmte Form verlangt. Zeilen mit \textbf{(E)} sind Ergebnispunkte für Ergebnisse direkt aus den Angaben: Das Ergebnis muss stimmen \emph{und} aus einem erkennbaren Weg hervorgehen. Zeilen mit «Folgezeile» bewerten Schritte aus eigenen Zwischenwerten: richtig weitergerechnet gibt den Punkt. Zeilen mit \textbf{(A)} sind Ablesepunkte: Der abgelesene Wert genügt. Alle übrigen Zeilen bewerten den Weg selbst.\par
4. Begründe jeden Abzug in einem Satz und nenne die Fehlerart (zum Beispiel «Minus vor der Klammer» oder «Fall $k = 2$ übersehen»). Schreib die Musterlösung nicht ab — zeig mir, wo mein Fehler liegt.\par
5. Gib zum Schluss eine Tabelle «Aufgabe | Punkte | Begründung», die Gesamtpunktzahl (von 23) und — nach der Selbsteinschätzung in Abschnitt 4 — welches Kapitel des Leitprogramms ich wiederholen soll. Keine Note.\par
6. Fehlt ein Foto oder ist eine Aufgabe unlesbar, sag es, statt sie zu bewerten.
\end{tcolorbox}

\needspace{16\baselineskip}
\section*{3 · Musterlösung und Punkteraster}
\textbf{Allgemein:} Ganze Punkte je Zeile. Ein Fehler wird nur einmal abgezogen (Folgefehler): Zeilen ohne Zeichen geben den
Punkt auch, wenn mit einem eigenen falschen Zwischenwert richtig weitergerechnet wird. \textbf{(E)}~= Ergebnispunkt für ein Ergebnis direkt aus den Angaben, nur wenn es stimmt und der Weg erkennbar ist; \textbf{(A)}~= Ablesepunkt. Derselbe Begriffsfehler in einer späteren Zeile derselben
Aufgabe gibt keinen zweiten Folgepunkt. Der ganze Test ist ohne Taschenrechner gestellt.
«Typische Fehler» nennt, welche Zeile 0 Punkte gibt, und die Punkte, die bleiben.
\textbf{Teilrichtige Zeilen:} Verlangt eine Zeile mehrere Angaben und stimmt nur ein Teil, gibt die Zeile 0 Punkte —
ausser die Zeile nennt ausdrücklich eine Aufteilung.

\begin{lpaufgabe}{G1}{Klammern auflösen und lösen}{3 P}
\begin{raster}
Klammern auflösen & 1 & $5x - 3x + 12 = 2x + 6 + x$, zusammengefasst $2x + 12 = 3x + 6$. Das Minus vor $3(x - 4)$ wirkt auf beide Glieder: $-3x + 12$.\\
Lösung & 1 & $\mid -2x - 6$: $6 = x$, $\mathbb{L} = \{6\}$. Folgezeile: Wer aus der eigenen Zeile richtig weiterrechnet, erhält den Punkt.\\
Probe & 1 & In der \emph{Ausgangsgleichung}: links $30 - 3 \cdot 2 = 24$, rechts $2 \cdot 9 + 6 = 24$. Gilt auch für eine richtig gerechnete Probe zur eigenen falschen Lösung, die den Fehler aufdeckt (beide Seiten verschieden, so benannt).\\
\end{raster}
\fehler{$-3(x - 4)$ als $-3x - 12$: $2x - 12 = 3x + 6$, $x = -18$ — erste Zeile 0, Lösungszeile als Folgezeile 1, Probe nach obiger Regel $\to$ höchstens 2 von 3\,P.
$-3(x - 4)$ als $-3x - 4$: $x = -10$ — ebenso.
Probe in einer umgeformten Zeile statt in der Ausgangsgleichung: Probezeile 0.}
\end{lpaufgabe}

\begin{lpaufgabe}{G2}{Gleichung mit Brüchen}{3 P}
\begin{raster}
Hauptnenner & 1 & Beide Seiten mal $12$, jedes Glied: $2(x - 2) = 3x - 12$ (der Zähler $x - 2$ als Ganzes). Gleichwertig: mal $24$ ($4(x - 2) = 6x - 24$).\\
Klammer, sortieren & 1 & $2x - 4 = 3x - 12$; $\mid -2x + 12$: $8 = x$\\
Lösung & 1 & $\mathbb{L} = \{8\}$ (Probe: $\tfrac{6}{6} = 1$ und $2 - 1 = 1$ — nicht verlangt). Folgezeile.\\
\end{raster}
\fehler{Die $1$ nicht mit $12$ multipliziert: $2(x - 2) = 3x - 1$ — Hauptnenner-Zeile 0; Klammerzeile ($x = -3$) und Lösungszeile als Folgezeilen je 1 $\to$ 2 von 3\,P.
Bruchstrich nicht als Klammer, $\tfrac{x - 2}{6} \cdot 12 = 2x - 2$: Hauptnenner-Zeile 0 (der Zähler muss als Ganzes mal 2); Folgezeilen ($x = 10$) je 1 $\to$ 2 von 3\,P.}
\end{lpaufgabe}

\begin{lpaufgabe}{G3}{Eine fremde Lösung prüfen}{3 P}
\begin{raster}
(a) Begründung & 1 & Die Umformung bis $0 = 0$ stimmt (links $2x - 2$, genau wie rechts). $0 = 0$ ist eine wahre Aussage, in der $x$ nicht mehr vorkommt: Sie gilt für \emph{jedes} $x$, nicht nur für $0$.\\
(a) Lösungsmenge & 1\E & $\mathbb{L} = \mathbb{R}$. Der Weg steht in der Aufgabe; das richtige Ergebnis genügt für diese Zeile, auch ohne Begründung.\\
(b) & 1 & Jeder Faktor $a \neq 2$, z.\,B. $3$: $2x - 2 = 3x - 2$ gibt $0 = x$, $\mathbb{L} = \{0\}$ — und das ist für jeden Faktor $a \neq 2$ dieselbe Lösung, weil $(2 - a)\,x = 0$. Faktor, Lösung $0$ und Rechnung nötig.\\
\end{raster}
\fehler{«$\{0\}$ ist richtig»: beide (a)-Zeilen 0.
Nur «$\mathbb{L} = \mathbb{R}$» ohne Begründung: Begründungszeile 0, Lösungszeile 1.
(b) Faktor $2$ behalten oder die Zahl $-2$ geändert: Zeile (b) 0. Richtiger Faktor, aber «keine Lösung» oder $\mathbb{L} = \{\,\}$ statt $\{0\}$: Zeile (b) 0.}
\end{lpaufgabe}

\begin{lpaufgabe}{G4}{Normalform und Wertetabelle}{4 P}
\begin{raster}
(a) Normalform & 1 & $\mid -9 + 2x$: $5x - 7 = 0$, $a = 5$, $b = -7$. Gleichwertig: jedes Vielfache mit derselben Lösung, z.\,B. $-5x + 7 = 0$ oder $10x - 14 = 0$, mit den passenden $a$, $b$. Form und $a$, $b$ nötig; $5x = 7$ ist keine Normalform (rechts muss $0$ stehen): Zeile 0.\\
(b) Lösung & 1 & $x = -\tfrac{b}{a} = \tfrac{7}{5} = 1.4$, $\mathbb{L} = \{1.4\}$. Folgezeile (auch aus $5x = 7$).\\
(c) Werte & 1 & $x = 1$: links $5$, rechts $7$, $5 < 7$. $x = 2$: links $8$, rechts $5$, $8 > 5$.\\
(c) Deutung & 1 & Das Vergleichszeichen kippt, also liegt die Lösung zwischen $1$ und $2$. Die Tabelle prüft nur die eingesetzten Werte; die Lösung $1.4$ ist keine ganze Zahl und kommt darin nicht vor — genau liefert sie nur die Rechnung.\\
\end{raster}
\fehler{(a) Vorzeichen beim Hinüberbringen: $5x + 7 = 0$ oder $x + 11 = 0$ — Zeile (a) 0; (b) mit dieser Form richtig weitergerechnet ($-1.4$ bzw. $-11$): Folgezeile 1.
(c) Nur «die Lösung liegt zwischen 1 und 2» ohne Grund, warum die Tabelle sie nicht genau trifft: Deutungszeile 0.}
\end{lpaufgabe}

\begin{lpaufgabe}{G5}{Zwei Tarife}{3 P}
\begin{raster}
(a) Ablesen & 1\A & Die Geraden schneiden sich zwischen $t = 3$ und $t = 4$ Stunden (bei etwa $3.3$).\\
(b) Rechnung & 1\E & $4t + 10 = 7t$, $10 = 3t$, $t = \tfrac{10}{3}$ h ($= 3$ h $20$ min); Preis $K = 7 \cdot \tfrac{10}{3} = \tfrac{70}{3} \approx 23.33$ CHF. Dauer \emph{und} Preis nötig; $3.33$ h und $23.33$ CHF zählen voll.\\
(c) & 1 & $K_A(5) = 30$ CHF, $K_B(5) = 35$ CHF: Tarif A ist günstiger. Mit Begründung (eingesetzt oder am Bild: rechts vom Schnittpunkt liegt A unten).\\
\end{raster}
\fehler{(b) $10 = 11t$ (Glieder addiert statt subtrahiert): Zeile (b) 0.
(b) Nur $t$ ohne Preis, oder $t = 3$ bzw. $t = 3.5$ gerundet: Zeile (b) 0.
(c) «A, weil …» ohne Vergleich der Werte oder des Bildes: Zeile (c) 0.}
\end{lpaufgabe}

\begin{lpaufgabe}{G6}{Lineare Gleichung mit Parameter}{4 P}
\begin{raster}
Sortieren & 1 & $\mid -2x - 4$: $k x - 2x = k^2 - 4$, also $(k - 2) \cdot x = k^2 - 4$.\\
Kritischer Wert & 1 & Der Faktor vor $x$ wird null bei $k = 2$. Dort: $0 \cdot x = 0$, wahr, $\mathbb{L} = \mathbb{R}$. Wert \emph{und} Lösungsmenge nötig; mit eigener rechter Seite richtig weitergedacht: Folgepunkt.\\
Faktorisieren & 1 & $k^2 - 4 = (k - 2)(k + 2)$ (dritte binomische Formel) — verlangt, weil die Aufgabe «so weit wie möglich vereinfachen» sagt.\\
Allgemeiner Fall & 1 & $k \neq 2$: $x = \tfrac{k^2 - 4}{k - 2}$, gekürzt $x = k + 2$, $\mathbb{L} = \{k + 2\}$. Folgezeile: Diese Zeile bewertet den Fall $k \neq 2$ samt Bedingung, gekürzt oder ungekürzt; das Kürzen bewertet die Zeile «Faktorisieren». Ohne die Bedingung $k \neq 2$ (oder «sonst» nach dem behandelten Fall $k = 2$): 0.\\
\end{raster}
\fehler{Ohne Fallunterscheidung durch $k - 2$ geteilt: Zeile «Kritischer Wert» 0 und Zeile «Allgemeiner Fall» 0 (Bedingung fehlt) $\to$ höchstens 2 von 4\,P.
Ungekürzt $\tfrac{k^2 - 4}{k - 2}$ mit Bedingung: Zeile «Faktorisieren» 0, «Allgemeiner Fall» 1.
Bei $k = 2$ «keine Lösung»: Zeile «Kritischer Wert» 0.
Sortierfehler $(k - 2) \cdot x = k^2 + 4$: Sortieren 0; $k = 2$: $0 = 8$, «keine Lösung» als Folgepunkt 1; Faktorisieren 0 ($k^2 + 4$ zerfällt nicht); Allgemeiner Fall $x = \tfrac{k^2 + 4}{k - 2}$ mit Bedingung als Folgezeile 1 $\to$ 2 von 4\,P.
$k = -2$ zusätzlich als Sonderfall mit «keine Lösung»: Zeile «Allgemeiner Fall» 0 — für $k = -2$ ist $x = 0$ eine Lösung.}
\end{lpaufgabe}

\begin{lpaufgabe}{G7}{Parameter bestimmen}{3 P}
\begin{raster}
(a) & 1 & $2k - 1 = 0$ bei $k = 0.5$; rechts bleibt $0.5 + 4 = 4.5 \neq 0$, also $0 \cdot x = 4.5$, keine Lösung. Wert \emph{und} Begründung nötig.\\
(b) Ansatz & 1 & $x = 3$ einsetzen: $(2k - 1) \cdot 3 = k + 4$, also $6k - 3 = k + 4$.\\
(b) Ergebnis & 1 & $5k = 7$, $k = \tfrac{7}{5} = 1.4$ (Probe: $1.8 \cdot 3 = 5.4 = 1.4 + 4$ — nicht verlangt). Folgezeile.\\
\end{raster}
\fehler{(a) $k = -0.5$ oder $k = 2$: Zeile (a) 0.
(b) Falsch eingesetzt, $2k - 1 = 3(k + 4)$: Ansatzzeile 0; daraus richtig $k = -13$: Ergebniszeile als Folgezeile 1.
(b) $(2k - 1) \cdot 3 = k + 4$ richtig, dann $6k - 1 = k + 4$ (Klammer): Ansatz 1, Ergebnis ($k = 1$) 0 — der Fehler liegt in dieser Zeile.}
\end{lpaufgabe}


\section*{4 · Selbsteinschätzung}
\begin{tabularx}{\linewidth}{@{}l X@{}}\toprule
20 – 23 P & Die geprüften Teile sitzen. Wo du Punkte verloren hast: das Kapitel dieser Aufgabe nochmals (Zuordnung unten).\\
15 – 19 P & Den schwächsten Teil nochmals: Tüfteln und Übungen des Kapitels, in dem du die meisten Punkte verloren hast.\\
10 – 14 P & Zurück zu den Kapiteln aller Aufgaben, in denen du Punkte verloren hast.\\
0 – 9 P & Zurück zu Kapitel 1 und von dort der Reihe nach weiter.\\\midrule
\multicolumn{2}{@{}p{\linewidth}@{}}{\small\textbf{Aufgabe $\to$ Kapitel:} G1 $\to$ Kapitel 1 und 2 (Umformen, Klammern); G2 $\to$ Kapitel 2 (Klammern und Brüche); G3 $\to$ Kapitel 3 (Lösungsfälle); G4, G5 $\to$ Kapitel 4 (Normalform, Tabelle, Graph); G6, G7 $\to$ Kapitel 5 (Parameterdiskussion)}\\\bottomrule
\end{tabularx}
\end{document}
